\chapter{Gap Identification and Problem Formulation}

\section{Gap Identified}
The review in Chapter 2 shows that the gap lies neither in modelling nor in language interpretation, both of which are mature, but in the connection between them. What is absent is a system in which a written requirement is interpreted by a language model, the resulting parameters drive a deterministic design calculation taken from published literature, and the calculation remains open to inspection by the user. Table~\ref{tab:tools} shows that none of the tools compared in Section~2.5 states all of these capabilities together.

Three deficiencies follow. First, parametric CAD regenerates geometry but verifies no failure mode. Second, systems that generate modelling code from a description produce dimensions that cannot be attributed to any formula. Third, where design calculation is automated, the working is generally hidden, so the result must be trusted rather than checked.

\section{Problem Formulation}
Given a design requirement expressed in ordinary English, the task is to produce a complete design of the specified machine joint such that every dimension is obtained from a published design formula or a named failure check, every applicable failure mode is verified against its allowable stress, the geometry is generated in a neutral exchange format, a first angle drawing sheet is produced, and the entire working is reported to the user in a form that permits independent verification.

\section{Conditions the Solution Must Satisfy}
\begin{enumerate}
  \item No dimension may be produced without an attributable basis, which is either a design formula with a cited source or an explicitly named failure check.
  \item The language model may interpret the requirement but may not select any dimension or perform any check.
  \item Any numerical value returned by the language model that does not appear in the original requirement must be discarded, so that an invented figure cannot enter the calculation.
  \item Every applicable failure mode must be evaluated on every run and its arithmetic reported, not merely its outcome.
  \item Where the empirical proportions used have no supporting source, the system must state that fact rather than imply a source it does not have.
\end{enumerate}
